\section{Demostración}

Procedemos a demostrar la optimalidad del algoritmo greedy propuesto, mediante inversiones. Para esto planteamos dos afirmaciones que una vez probadas, nos permitirán concluir el resultado. 

Decimos que una secuencia de videos tiene una inversión, si dentro de la misma tenemos dos videos $V_i, V_j$ tal que $i < j$ pero $a_i < a_j$.
\\

\textbf{Comparación de secuencias sin inversiones}

\textbf{Postulado: } Toda secuencia sin inversiones tiene el mismo tiempo de finalización.

\textbf{Demostración: } Por no tener inversiones, solo pueden variar la posición de videos que tienen el mismo $a_i$, y van a tener que ser consecutivos. El ultimo de los videos es el que tiene mayor $a_i$ entre ellos.

\textbf{Ejemplo}

Dados 3 videos ($v_1, v_2$ y $v_3$), para los cuales: 

$a_1 = a_2 = a_3 = 4$

$s_1 = 1, s_2 = 2$ y $s_3 = 3$

$v_1, v_2, v_3$ \rightarrow C = 10

$v_2, v_1, v_3$ \rightarrow C = 10

$v_1, v_3, v_2$ \rightarrow C = 10

(...) 
el tiempo va a ser el mismo con todas sus premutaciones.
\\

\textbf{Secuencia sin inversiones optima}

Sea una solución cualquiera que no respeta el orden decreciente de $a_i$. Entonces, existe al menos un par de videos consecutivos $i, j$ tal que $i$ aparece antes que $j$ y:
\[
a_i < a_j.
\]

Sea $T$ el tiempo acumulado de Scaloni antes de procesar $i$ y $j$.

\textbf{Orden original:} $i \rightarrow j$
\[
C_i = T + s_i + a_i, \quad
C_j = T + s_i + s_j + a_j.
\]
El tiempo máximo en este bloque es:
\[
\max(C_i, C_j) = \max(T + s_i + a_i,\; T + s_i + s_j + a_j).
\]

\textbf{Orden intercambiado:} $j \rightarrow i$
\[
C'_j = T + s_j + a_j, \quad
C'_i = T + s_j + s_i + a_i.
\]
El tiempo máximo pasa a ser:
\[
\max(C'_j, C'_i) = \max(T + s_j + a_j,\; T + s_j + s_i + a_i).
\]

Como $a_j > a_i$, se tiene que:
\[
T + s_i + s_j + a_j \ge T + s_i + a_i,
\]
por lo que en el orden original domina $C_j$.

Además,
\[
T + s_j + a_j \le T + s_i + s_j + a_j,
\]
y
\[
T + s_j + s_i + a_i \le T + s_i + s_j + a_j,
\]
por lo tanto:
\[
\max(C'_j, C'_i) \le \max(C_i, C_j).
\]

Esto demuestra que intercambiar $i$ y $j$ no aumenta el tiempo máximo.

Repitiendo este proceso para eliminar todas las inversiones, se obtiene un orden en el cual los $a_i$ están en orden decreciente, sin aumentar el tiempo total.

\textbf{Conclusión:} Existe una solución óptima en la que los videos están ordenados en orden decreciente según $a_i$, lo que justifica la estrategia greedy.
